\documentclass[10pt,a4paper]{amsart}
\usepackage{amsmath,amssymb,amsthm,graphicx}
\usepackage{verbatim}
\newenvironment{reference}{\comment}{\endcomment}
%\newenvironment{reference}{}{}
\newenvironment{slogan}{\comment}{\endcomment}
\newenvironment{history}{\comment}{\endcomment}

% For commutative diagrams we use Xy-pic
\usepackage[all]{xy}

% We use 2cell for 2-commutative diagrams.
\xyoption{2cell}
% Theorem environments.
%
\theoremstyle{plain}
\newtheorem{theorem}[subsection]{Theorem}
\newtheorem{proposition}[subsection]{Proposition}
\newtheorem{lemma}[subsection]{Lemma}

\theoremstyle{definition}
\newtheorem{definition}[subsection]{Definition}
\newtheorem{example}[subsection]{Example}
\newtheorem{exercise}[subsection]{Exercise}
\newtheorem{situation}[subsection]{Situation}

\theoremstyle{remark}
\newtheorem{remark}[subsection]{Remark}
\newtheorem{remarks}[subsection]{Remarks}

\numberwithin{equation}{subsection}

% Macros
%
\def\lim{\mathop{\mathrm{lim}}\nolimits}
\def\colim{\mathop{\mathrm{colim}}\nolimits}
\def\Spec{\mathop{\mathrm{Spec}}}
\def\Hom{\mathop{\mathrm{Hom}}\nolimits}
\def\Ext{\mathop{\mathrm{Ext}}\nolimits}
\def\SheafHom{\mathop{\mathcal{H}\!\mathit{om}}\nolimits}
\def\SheafExt{\mathop{\mathcal{E}\!\mathit{xt}}\nolimits}
\def\Sch{\mathit{Sch}}
\def\Mor{\mathop{\mathrm{Mor}}\nolimits}
\def\Ob{\mathop{\mathrm{Ob}}\nolimits}
\def\Sh{\mathop{\mathit{Sh}}\nolimits}
\def\NL{\mathop{N\!L}\nolimits}
\def\CH{\mathop{\mathrm{CH}}\nolimits}
\def\proetale{{pro\text{-}\acute{e}tale}}
\def\etale{{\acute{e}tale}}
\def\QCoh{\mathit{QCoh}}
\def\Ker{\mathop{\mathrm{Ker}}}
\def\Im{\mathop{\mathrm{Im}}}
\def\Coker{\mathop{\mathrm{Coker}}}
\def\Coim{\mathop{\mathrm{Coim}}}

% Boxtimes
%
\DeclareMathSymbol{\boxtimes}{\mathbin}{AMSa}{"02}

%
% Macros for moduli stacks/spaces
%
\def\QCohstack{\mathcal{QC}\!\mathit{oh}}
\def\Cohstack{\mathcal{C}\!\mathit{oh}}
\def\Spacesstack{\mathcal{S}\!\mathit{paces}}
\def\Quotfunctor{\mathrm{Quot}}
\def\Hilbfunctor{\mathrm{Hilb}}
\def\Curvesstack{\mathcal{C}\!\mathit{urves}}
\def\Polarizedstack{\mathcal{P}\!\mathit{olarized}}
\def\Complexesstack{\mathcal{C}\!\mathit{omplexes}}
% \Pic is the operator that assigns to X its picard group, usage \Pic(X)
% \Picardstack_{X/B} denotes the Picard stack of X over B
% \Picardfunctor_{X/B} denotes the Picard functor of X over B
\def\Pic{\mathop{\mathrm{Pic}}\nolimits}
\def\Picardstack{\mathcal{P}\!\mathit{ic}}
\def\Picardfunctor{\mathrm{Pic}}
\def\Deformationcategory{\mathcal{D}\!\mathit{ef}}
\usepackage{cleveref}
\crefname{originallemma}{original lemma}{original lemmas}
\Crefname{originallemma}{Original lemma}{Original lemmas}
\makeatletter
\newlabel{lemma-codimension}{{00P2}{source}{}{sourcepointer-1}{}}
\newlabel{lemma-codimension@cref}{{[originallemma][0][]00P2}{[1][0][]source}{}{}{}}
\makeatother
\crefname{originallemma}{original lemma}{original lemmas}
\Crefname{originallemma}{Original lemma}{Original lemmas}
\makeatletter
\newlabel{lemma-dimension-base-fibre-equals-total}{{00ON}{source}{}{sourcepointer-2}{}}
\newlabel{lemma-dimension-base-fibre-equals-total@cref}{{[originallemma][0][]00ON}{[1][0][]source}{}{}{}}
\makeatother
\title{The local dimension at a point is unchanged by extending the base field}
\author{The Stacks Project Authors}
\date{stacks-project:algebra.tex:local-snapshot-20261012; GNU FDL 1.2 or later}
\begin{document}
\maketitle
\noindent\textit{Editorial note.} All delivered mathematical text is the author's own, copied from the pinned Stacks Project chapter file algebra.tex; no statement or proof step was rewritten, reordered or removed.\par\medskip
\noindent\textit{Editorial note.} Preamble provenance, as required for this source: the frozen snapshot stages the chapter's own preamble file beside it (preamble.tex, 217 lines, 8116 bytes, the same shared Stacks-wide preamble the published chow and etale-cohomology cards use); the first line of the pinned chapter file is \textbackslash{}input\{preamble\}, so no preamble text lives inside the chapter itself. The card takes that file's lines 8-18 (the verbatim package and the three \textbackslash{}comment-based environments reference/slogan/history at lines 9-12, section 13.2, together with the Xy-pic setup) and its lines 154 to the runtime end of the file (the native theorem declarations and mathematical macros). Not one character was changed, no package was added, no macro was defined, and the excerpt serves only as preamble, never as a span coordinate.\par\medskip
\noindent\textit{Editorial note.} The same-chapter pointers are located at their \textbackslash{}label lines inside the pinned chapter file; the cross-chapter pointers are located in the project's own tag list staged in the snapshot (sources/stacks-chow/tags.txt), at the verbatim lines carrying their canonical tags.\par\medskip
\noindent\textit{Editorial note.} The counted proof refers to same-chapter statements by label; they are recorded as navigation stubs pointed at their \textbackslash{}label lines, and the arguments using them stay inside the counted span.\par\medskip
\noindent\textit{Editorial note.} The counted spans contain no \textbackslash{}cite, so the card carries no bibliography.\par\medskip
\noindent\textit{Allowed standard prerequisites (editorial).}\begin{itemize}
\item Presentation of the base change by extending the ideal. Used at the start of the counted proof (lines 28537-28540). Describes the extended algebra as a quotient of a polynomial ring by the extended ideal.
\item The dimension formula for a point of a finite type algebra over a field. Used at the end of the counted proof (lines 28560-28571). Expresses both local dimensions through the same numerical data, which is what makes them equal.
\end{itemize}\medskip
\ifdefined\hypertarget\hypertarget{sourcepointer-1}{}\fi
\noindent\textit{Source reference.} Original lemma 00P2: sources/stacks-algebra/pinned-original-source.tex, source lines 28487--28487. Navigation pointer for \textbackslash{}ref\{lemma-codimension\} used inside the counted span (line 28563); the argument that consumes it is retained in the span.\par\medskip
\ifdefined\hypertarget\hypertarget{sourcepointer-2}{}\fi
\noindent\textit{Source reference.} Original lemma 00ON: sources/stacks-algebra/pinned-original-source.tex, source lines 27638--27638. Navigation pointer for \textbackslash{}ref\{lemma-dimension-base-fibre-equals-total\} used inside the counted span (line 28557); the argument that consumes it is retained in the span.\par\medskip
\begin{lemma}
\label{lemma-dimension-at-a-point-preserved-field-extension}
Let $k$ be a field.
Let $S$ be a finite type $k$-algebra.
Set $X = \Spec(S)$.
Let $K/k$ be a field extension.
Set $S_K = K \otimes_k S$, and $X_K = \Spec(S_K)$.
Let $\mathfrak q \subset S$ be a prime corresponding to $x \in X$
and let $\mathfrak q_K \subset S_K$ be a prime corresponding
to $x_K \in X_K$ lying over $\mathfrak q$.
Then $\dim_x X = \dim_{x_K} X_K$.
\end{lemma}

\begin{proof}
Choose a presentation $S = k[x_1, \ldots, x_n]/I$.
This gives a presentation
$K \otimes_k S = K[x_1, \ldots, x_n]/(K \otimes_k I)$.
Let $\mathfrak q_K' \subset K[x_1, \ldots, x_n]$,
resp.\ $\mathfrak q' \subset k[x_1, \ldots, x_n]$ be
the corresponding primes. Consider the following
commutative diagram of Noetherian local rings
$$
\xymatrix{
K[x_1, \ldots, x_n]_{\mathfrak q_K'} \ar[r] &
(K \otimes_k S)_{\mathfrak q_K} \\
k[x_1, \ldots, x_n]_{\mathfrak q'} \ar[r] \ar[u] &
S_{\mathfrak q} \ar[u]
}
$$
Both vertical arrows are flat because they are localizations of
the flat ring maps $S \to S_K$ and
$k[x_1, \ldots, x_n] \to K[x_1, \ldots, x_n]$.
Moreover, the vertical arrows have the same fibre rings.
Hence, we see from
Lemma \ref{lemma-dimension-base-fibre-equals-total} that
$\text{height}(\mathfrak q') - \text{height}(\mathfrak q)
= \text{height}(\mathfrak q_K') - \text{height}(\mathfrak q_K)$.
Denote $x' \in X' = \Spec(k[x_1, \ldots, x_n])$
and $x'_K \in X'_K = \Spec(K[x_1, \ldots, x_n])$
the points corresponding to $\mathfrak q'$ and
$\mathfrak q_K'$. By Lemma \ref{lemma-codimension} and what we showed
above we have
\begin{eqnarray*}
n - \dim_x X & = & \dim_{x'} X' - \dim_x X \\
& = & \text{height}(\mathfrak q') - \text{height}(\mathfrak q) \\
& = & \text{height}(\mathfrak q_K') - \text{height}(\mathfrak q_K) \\
& = & \dim_{x'_K} X'_K - \dim_{x_K} X_K \\
& = & n - \dim_{x_K} X_K
\end{eqnarray*}
and the lemma follows.
\end{proof}

\end{document}
